\section{Physical reservoir and accuracy criterion}
\label{sec:framework}

\subsection{Canonical reference and the exact composite marginal}

Let $(\Omega_N,\mathcal F_N,\mu_N)$ describe the subsystem states, with
measurable finite energy $E_N:\Omega_N\to\R$. The reference measure may be
counting measure or a classical phase-space measure. At the specified
inverse temperature $\beta_N>0$, assume
\begin{equation}
 P_N(\dd x)=\frac{e^{-\beta_NE_N(x)}}{\mathcal Z_N}\mu_N(\dd x),
 \qquad 0<\mathcal Z_N<\infty.
 \label{eq:canonical}
\end{equation}
Throughout, $\beta_N\to\beta\in(0,\infty)$; a fixed temperature is included.
Allowing this specified sequence is useful for finite-size equal-weight
coexistence. The target law is chosen before the reservoir is tuned.

We neglect the interaction energy between subsystem and reservoir. The
reservoir has the exact surface density of states
\begin{equation}
 \omega_{B,N}(U)=A_N U^{c_N}\ind_{\{U>0\}},
 \qquad A_N>0,\quad c_N>0.
 \label{eq:bath-dos}
\end{equation}
Here ``surface'' means the density of states at an energy, rather than its
integral up to that energy. Integrating a microcanonical energy shell for
the composite gives, at composite energy $\mathcal E_N$, the subsystem law
\begin{equation}
 Q_N(\dd x)=\frac{L_N(E_N(x))}{Z_N}P_N(\dd x),\qquad
 L_N(e)=
 \begin{cases}
 e^{\beta_N e}(\mathcal E_N-e)^{c_N},&e<\mathcal E_N,\\
 0,&e\ge\mathcal E_N,
 \end{cases}
 \quad Z_N=\E_{P_N}L_N(E_N).
 \label{eq:exact-marginal}
\end{equation}
The bath exponent $c_N$ is prescribed. Only $\mathcal E_N$ is adjustable.
For every finite $\mathcal E_N$, the function $L_N$ is bounded: it is
continuous below the cutoff and tends to zero both at the cutoff and as
$e\to-\infty$. Thus $Z_N<\infty$ automatically. The total energy is
admissible when $P_N(E_N<\mathcal E_N)>0$, which is exactly the condition
$Z_N>0$.

To specify the heat-capacity convention, write
$S_{B,N}(U)=k_B\log\omega_{B,N}(U)$, with a fixed reference energy absorbed
in the additive constant. Its temperature $T_B$ satisfies
$1/T_B=\partial_U S_{B,N}=k_Bc_N/U$. Hence
\begin{equation}
 \beta_B=\frac{1}{k_BT_B}=\frac{c_N}{U},\qquad
 C_B=\frac{\dd U}{\dd T_B}=k_Bc_N.
 \label{eq:capacity-convention}
\end{equation}
For $f_N>2$ independent classical quadratic degrees of freedom,
$c_N=f_N/2-1$. At a canonical inverse temperature $\beta_{\mathrm{can}}$,
the same reservoir has an energy density proportional to
$U^{c_N}e^{-\beta_{\mathrm{can}}U}\ind_{\{U>0\}}$, with mean energy
$(c_N+1)/\beta_{\mathrm{can}}$ and canonical heat capacity $k_B(c_N+1)$. This finite
offset follows from the chosen surface convention and does not affect the
diverging scales proved below. The theorems concern the specified density
of states and exponent; they do not require a choice between the
competing surface and volume definitions of microcanonical entropy
compared by Hilbert, H\"anggi, and Dunkel
\cite{HilbertHanggiDunkel2014,DunkelHilbert2014}. No quadratic expansion of
the bath entropy is assumed in~\eqref{eq:exact-marginal}.

\subsection{Full-state and energy-law accuracy coincide}

For probability measures $P,Q$ on the same space, use
\begin{equation}
 \TV(P,Q)=\sup_{A}|P(A)-Q(A)|
 =\frac12\int|\dd P-\dd Q|.
 \label{eq:tv-definition}
\end{equation}
The last integral uses any common dominating measure. Let
$P_N^E,Q_N^E$ be the pushforwards by $E_N$.

\begin{lemma}[Energy is sufficient for reservoir reweighting]
\label{lem:energy-tv}
For the exact marginal~\eqref{eq:exact-marginal},
\begin{equation}
 \TV(Q_N,P_N)=\TV(Q_N^E,P_N^E)
 =\frac12\E_{P_N}\left|\frac{L_N(E_N)}{Z_N}-1\right|.
 \label{eq:energy-tv}
\end{equation}
\end{lemma}
\begin{proof}
The Radon--Nikodym derivative on the state space is $L_N(E_N)/Z_N$.
Its pushforward derivative is the same function of energy: for every
bounded measurable $f$,
\[
 \E_{Q_N}f(E_N)=\frac1{Z_N}\int f(e)L_N(e)P_N^E(\dd e).
\]
The two $L^1$ formulas for total variation are therefore identical.
\end{proof}

We consequently work with energy laws without losing any information
about this full-state distance. Matching a mean energy, a pressure,
or a few local observables need not determine this distance. Conversely,
total variation controls uniformly bounded observables; convergence of
expectations of unbounded observables requires additional tail control,
such as uniform integrability under both laws. The equality in
Lemma~\ref{lem:energy-tv} depends on the energy-only likelihood and need
not hold when coupling changes the subsystem Hamiltonian.

\begin{lemma}[Normalization criterion]
\label{lem:normalization}
Suppose $W_N\ge0$ and $\E_{P_N}|W_N-1|\to0$. Then
$\E_{P_N}W_N\to1$, and the probability law with derivative
$W_N/\E_{P_N}W_N$ converges to $P_N$ in total variation. Conversely,
total-variation convergence of laws with derivatives $R_N$ implies
$R_N\to1$ in $P_N$-probability. In particular, if $0\le W_N\le1$
and $W_N\to1$ in probability, the first assertion applies.
\end{lemma}
\begin{proof}
Writing $z_N=\E W_N$, we have $|z_N-1|\le\E|W_N-1|$ and
\[
 \frac12\E\left|\frac{W_N}{z_N}-1\right|
 \le\frac{\E|W_N-1|+|z_N-1|}{2z_N}.
\]
The converse follows from Markov's inequality and~\eqref{eq:tv-definition}.
The last assertion follows from boundedness and convergence in probability.
\end{proof}

All asymptotic scales refer to fixed energy units. We write $u_N\gg v_N$
when $u_N/v_N\to\infty$, and $u_N\asymp v_N$ when their ratio is bounded
above and below by positive constants. In dimensional form, expressions
such as $b_N^2/c_N$ or $\Delta_Ns_N/c_N$ are multiplied by $\beta_N^2$.
Because $\beta_N$ has a positive finite limit, these factors do not change
any asymptotic condition.

\subsection{Two exact choices of composite energy}

For one relevant energy center $a_N$, choose
\begin{equation}
 \mathcal E_N=a_N+\frac{c_N}{\beta_N}.
 \label{eq:centered-calibration}
\end{equation}
After division by $L_N(a_N)$, the weight becomes
\begin{equation}
 W_N(e)=
 \begin{cases}
 \exp[\beta_N(e-a_N)]
 \left(1-\dfrac{\beta_N(e-a_N)}{c_N}\right)^{c_N},
 &\beta_N(e-a_N)<c_N,\\
 0,&\beta_N(e-a_N)\ge c_N.
 \end{cases}
 \label{eq:centered-weight}
\end{equation}
Since $\log(1-y)\le-y$ for every $y<1$,
\begin{equation}
 0\le W_N(e)\le1\qquad(e\in\R).
 \label{eq:global-centered-bound}
\end{equation}
This global bound, including arbitrarily low energies, is what permits
a theorem without moment or tail assumptions.

For two centers $e_{-,N}<e_{+,N}$ with gap
$\Delta_N=e_{+,N}-e_{-,N}$, use instead
\begin{equation}
 \mathcal E_N=e_{-,N}+
 \frac{\Delta_N}{1-e^{-\beta_N\Delta_N/c_N}}
 =e_{+,N}+\frac{\Delta_N}{e^{\beta_N\Delta_N/c_N}-1}.
 \label{eq:secant-calibration}
\end{equation}
Then $L_N(e_{-,N})=L_N(e_{+,N})>0$. Define the reference-normalized residual
log-weight, obtained by subtracting its value at $e_{-,N}$, on its feasible half-line by
\begin{equation}
 h_N(e)=\beta_N(e-e_{-,N})+
 c_N\log\frac{\mathcal E_N-e}{\mathcal E_N-e_{-,N}},
 \quad e<\mathcal E_N,
 \label{eq:secant-h}
\end{equation}
and set $e^{h_N(e)}=0$ at and above the cutoff. This reference normalization
is distinct from probability normalization by $\E_{P_N}e^{h_N(E_N)}$.
It satisfies
\begin{equation}
 h_N(e_{-,N})=h_N(e_{+,N})=0,\qquad
 h_N'(e)=\beta_N-\frac{c_N}{\mathcal E_N-e},\qquad
 h_N''(e)=-\frac{c_N}{(\mathcal E_N-e)^2}.
 \label{eq:secant-derivatives}
\end{equation}
In particular, $h_N\ge0$ between the centers and $h_N\le0$ outside
them whenever feasible. Balancing the two center values is a physical
choice of total energy, not an independently imposed bias. It need not
balance the integrated probabilities of two broad phases.

\begin{lemma}[Bounds for the two-center calibration]
\label{lem:secant-bounds}
Suppose $\Delta_N\to\infty$ and $c_N/\Delta_N\to\infty$. For all large
$N$, with a constant $C$ depending only on fixed bounds for $\beta_N$,
\begin{align}
 |h_N'(e_{\pm,N})|&\le C\Delta_N/c_N,
 \label{eq:endpoint-slope-bound}\\
 \sup_{e<\mathcal E_N}h_N(e)&\le C\Delta_N^2/c_N,
 \label{eq:global-amplification}\\
 h_N(e)&\le C\frac{\Delta_N}{c_N}|e-e_{i,N}|
 \quad(i\in\{-,+\},\ e<\mathcal E_N).
 \label{eq:tangent-bound}
\end{align}
If $s_N=o(\Delta_N)$ and $c_N\gg\Delta_Ns_N$, then for each fixed $R$,
\begin{equation}
 \sup_{|e-e_{i,N}|\le Rs_N}|h_N(e)|\longrightarrow0,
 \label{eq:phase-flatness}
\end{equation}
provided $s_N\to\infty$. These windows are below the cutoff for all
sufficiently large $N$.
\end{lemma}
\begin{proof}
Put $t_N=\beta_N\Delta_N/c_N\to0$. The residual reservoir energies at
the lower and upper centers are $\Delta_N/(1-e^{-t_N})$ and
$\Delta_N/(e^{t_N}-1)$, respectively. Thus their inverse temperatures
are
\[
 \beta_N\frac{1-e^{-t_N}}{t_N}
 \quad\hbox{and}\quad
 \beta_N\frac{e^{t_N}-1}{t_N}.
\]
Their difference from $\beta_N$ is $O(t_N)$, giving
\eqref{eq:endpoint-slope-bound}. On the interval between the centers,
the residual reservoir energy is $(c_N/\beta_N)[1+O(t_N)]$, so
$0\le-h_N''\le C/c_N$. Interpolation between two zero endpoints gives
\[
 0\le h_N(e)\le\frac{C}{2c_N}
 (e-e_{-,N})(e_{+,N}-e).
\]
For completeness, subtract the right side from $h_N$: the difference
is convex and vanishes at both endpoints, hence is nonpositive throughout
the interval.
Outside the interval, concavity gives $h_N\le0$. This proves
\eqref{eq:global-amplification}. A concave differentiable function lies
below each tangent, proving~\eqref{eq:tangent-bound}.

The cutoff is at distance asymptotic to $c_N/\beta_N$ from either
center. On every indicated window, $-h_N''\le C/c_N$ for large $N$.
Taylor's theorem and the endpoint slope estimate give
\[
 |h_N(e)|\le C\frac{\Delta_NRs_N+R^2s_N^2}{c_N}\longrightarrow0.
\]
\end{proof}
